\section{Materials and Methods}

\subsection{Selective execution}

Selective execution is the central design principle of VariantFlow. A VCF record
stores every field for every sample on a single line, so a conventional parser
must walk the whole line before it can return any part of it. Most analyses,
however, depend on a small fraction of that line: a quality filter needs
\texttt{QUAL}, a missingness summary needs only the genotype of each sample, and
an allele-frequency calculation needs neither read depths nor genotype
qualities. The remaining bytes are decoded and discarded.

Each query, whether a filter expression, a population-genetic statistic, or a
column projection, is therefore compiled into a typed predicate tree that
identifies the specific fields on which the result depends (Figure~1).
VariantFlow then streams once through the file and, for each record, decodes only
those fields. Rather than materialising the whole record as strings, it treats
each record as a borrowed byte-slice view and locates the required site, INFO,
and FORMAT fields by scanning the raw bytes in place, so no allocation is made
for a field that is never read. Because the skipped fields cannot, by
construction, affect the result, the output is identical to that of a tool that
decodes every field. That equivalence is a design property rather than an
approximation: VariantFlow reads fewer fields, never fewer records, and never
samples or estimates. It is verified empirically in the Results.

\subsection{Field-selective parsing of per-sample data}

The cost that selective execution avoids grows with cohort size, because the
per-sample block is the part of a VCF record that scales. In a 3\,202-sample
cohort each record carries thousands of \texttt{GT:DP:GQ:AD} entries, and a
filter on read depth alone must otherwise traverse all of them. VariantFlow
parses the FORMAT schema once per record, caches the position of the requested
key such as \texttt{DP}, and reads that value by offset across all samples,
avoiding a repeated scan of the layout for every sample. In its line-preserving
mode, records that pass a filter are written in input order with their original
text unchanged, so a filtered file remains byte-comparable with its source.

\subsection{Compressed input, parallelism, and query-aware indexing}

Compressed input is read through a native block-gzip (BGZF) reader with a
worker-thread count capped automatically to the host. CPU-bound predicates are
evaluated in ordered batches so that output order is preserved regardless of the
number of threads.

For filters that are applied repeatedly and reject most records, an optional
\texttt{.vfi} sidecar index stores per-block metadata: the genomic range,
\texttt{QUAL} bounds, the set of \texttt{FILTER} values present, and summaries of
selected INFO fields. The query planner consults this metadata and skips an
entire compressed block whenever it proves that no record within it can satisfy
the filter. The distinction from a coordinate index such as tabix
\citep{tabix} is that the skipping criterion need not be positional. A tabix or
CSI index answers ``which blocks overlap this region''; it cannot answer ``which
blocks contain a record whose \texttt{FILTER} is \texttt{PASS}'' or ``whose
allele frequency exceeds a threshold'', because those properties are not
functions of genomic position. Storing per-block summaries of exactly the fields
that predicates reference is what makes non-positional skipping possible, and it
is the reason the largest speedups in this study arise for predicates that
exclude most of the file. When few blocks can be excluded, the planner falls back
to a sequential scan.

Through an optional HTSlib backend \citep{htslib}, VariantFlow additionally reads
BCF (the binary encoding of VCF), restricts reading to a single indexed region,
and writes BGZF output.

\subsection{Population-genetic statistics and columnar export}

For population-genetic analysis, the VCFtools-compatible subcommands operate on
compact diploid biallelic genotype summaries and reproduce VCFtools output for
allele frequency, missingness, heterozygosity, Hardy--Weinberg equilibrium,
nucleotide diversity, Tajima's $D$, linkage disequilibrium, and $F_{ST}$. Each is
computed in a single streaming pass over sites or windows, so peak memory is set
by the number of samples rather than the number of records. For repeated
analytical queries, VariantFlow exports typed columns in the Parquet format
\citep{parquet}, which columnar engines such as DuckDB \citep{duckdb} can query
efficiently without re-parsing the VCF.

\subsection{Interface and usage}

The command-line interface follows the compositional, single-purpose style common
in genomics, exposing the operations as a small tree of filtering, indexing,
input/output, and population-genetic subcommands (Figure~2).

A representative FORMAT-aware filter keeps sites that pass a quality threshold
and have at least one sample sequenced to sufficient depth:

{\small
\texttt{variantflow filter chr22.vcf.gz \textbackslash}\\
\texttt{~~--where "QUAL > 30 \&\& ANY(FORMAT/DP > 20)" -o chr22.filtered.vcf}
}

\noindent
The output is a VCF whose passing records are byte-identical to the
corresponding input lines. A two-population $F_{ST}$ computation takes two files
of sample identifiers and writes a tab-separated table:

{\small
\texttt{variantflow fst chr22.vcf.gz --pop AFR.txt --pop EUR.txt -o fst.tsv}
}

\noindent
Each row of \texttt{fst.tsv} gives the chromosome, position, and the
Weir--Cockerham estimate for that site, so a value of, for example, $0.31$ at a
given position indicates that roughly a third of the genetic variance there is
partitioned between the two populations rather than within them. The same
invocation pattern applies to the other summaries, which differ only in the
statistic named and the options it requires. A worked end-to-end analysis of
1000~Genomes chromosome~22, from installation through interpretation, is provided
as a supplementary user guide.

VariantFlow is written in Rust and builds from source with Cargo on Linux and
macOS; prebuilt binaries and a container image are distributed with each release.
The optional HTSlib backend is enabled at compile time and is not required for
any result reported here.

\subsection{Benchmark and validation protocol}

All comparisons reported here are drawn from benchmark reports tracked in the
repository, each recording the exact commands, the competitor version, and the
correctness check applied. A speedup is reported only for an operation whose
output has first been verified against a reference tool (Supplementary
Table~S1); operations that trail a competitor are treated as optimisation
targets rather than claims, and are reported as measured.

Population-genetic benchmarks were run on an Apple M5 Max workstation and the
index-assisted and columnar benchmarks on a Docker/Linux host. Within every
comparison both tools were run on the same host, so no ratio spans machines. The
full system configuration, tool versions, and dataset shapes are given in
Supplementary Table~S2. Timing is single-threaded wall-clock; memory is peak
resident set size, and correctness was assessed by normalising both outputs and
comparing them field by field.
